% TOP_PROBLEM: {"id": 6700007, "title": "Scalar Curvature Question [?7]: But deeper structures (if they exist at all) that lie at the roots of Dirac operators and of minimal hypersurf", "category_id": 6, "category": {"id": 6, "name": "geometry", "display_name": "Geometry", "description": "Euclidean and non-Euclidean geometry, geometric structures.", "slug": "geometry", "order_index": 6, "created_at": "2026-07-31T15:26:25.671Z"}, "status": "open"}
% FIRST_POSTED: null
% ENTRY_KIND: statement_only
% Imported from: batch14.tex; UnsolvedMath ID 6700007
\documentclass[11pt]{article}
\usepackage{fontspec}
\setmainfont{Latin Modern Roman}
\setsansfont{Latin Modern Sans}
\usepackage{amsmath,amssymb,mathtools}
\usepackage{unicode-math}
\setmathfont{Latin Modern Math}
\usepackage{xurl}
\usepackage[hidelinks]{hyperref}
\newcommand{\sref}[2]{\hyperlink{source:#1:#2}{[S#2]}}
\newcommand{\cataloguescope}[1]{\par\noindent\textbf{Question.} #1\par}
\begin{document}
% UnsolvedMath ID 6700007; source catalogue code AMR-066-0007
\setcounter{section}{6700006}
\section{Scalar Curvature Question [?7]: But deeper structures (if they exist at all) that lie at the roots of Dirac operators and of minimal hypersurf}\label{6700007}
{\small\sffamily UnsolvedMath ID 6700007 · Geometry}
\subsection{Definitions and mathematical statement}

\paragraph{Shared notation.}$\mathbb N=\{1,2,\ldots\}$, and $\mathbb Z,\mathbb Q,\mathbb R,\mathbb C$ denote the integers, rationals, reals and complex numbers. Any other conventions are specified locally.


A Dirac operator is a first-order elliptic operator on spinor fields, or on spinors twisted by an auxiliary bundle, used to extract scalar-curvature constraints through its index and square. A minimal hypersurface is a submanifold of codimension one with vanishing mean curvature, equivalently stationary area under local variations. Minimal varieties in the source allow the corresponding higher-dimensional and singular geometric-measure-theory setting. Both techniques supply restrictions on scalar curvature, but the question does not prescribe a particular common structure or an equivalence theorem.
\paragraph{Statement recorded in the supplied source.}
Identify deeper structures, if any, underlying both Dirac-operator methods and minimal-hypersurface methods in scalar curvature. Is there a direct link between Dirac operators and minimal varieties, or is their joint appearance in this setting accidental? \sref{6700007}{2}
\subsection{Short English statement}
Seek a structural explanation, if one exists, for why Dirac operators and minimal hypersurfaces both govern scalar-curvature phenomena.
\subsection{Sources}
\begin{description}
\item[\hypertarget{source:6700007:1}{S1}] ULAM.AI and the UnsolvedMath Contributors, UnsolvedMath: A Curated Collection of Open Mathematics Problems, record 6700007 (AMR-066-0007). CC BY 4.0. \url{https://huggingface.co/datasets/ulamai/UnsolvedMath}.
\item[\hypertarget{source:6700007:2}{S2}] Misha Gromov, Four Lectures on Scalar Curvature (2021 version). Section 3.20, Question 2 and its surrounding discussion, p. 204. \url{https://arxiv.org/pdf/1908.10612v6}.
\end{description}
\label{end:6700007}
\end{document}
